\chapter{Microkernel dan Plugin}
\label{bab:plugin}

\section*{Tujuan Pembelajaran}

\begin{enumerate}
  \item Menjelaskan pembagian antara inti dan \textit{plugin} beserta kontrak
    yang menghubungkan keduanya.
  \item Membedakan aktivasi awal dari aktivasi lambat, lalu menentukan kapan
    masing-masing bentuk layak dipilih.
  \item Mengukur penyebutan \textit{plugin} di dalam inti, akibat satu
    \textit{plugin} yang rusak, dan selisih waktu siap antara kedua cara
    aktivasi.
\end{enumerate}

\section{Pendahuluan}

Microkernel memisahkan sebuah sistem menjadi inti yang kecil beserta sejumlah
bagian tambahan yang dapat dipasang dan dicabut. Istilah tersebut lahir di
dunia sistem operasi, yaitu ketika penjadwalan dan pengelolaan memori
dipertahankan di dalam kernel, sedangkan sistem berkas dan pengandar perangkat
dipindahkan ke luar sebagai proses tersendiri.

Bab ini hanya menyinggung sistem operasi sekilas, sebab bentuk yang jauh lebih
sering ditemui mahasiswa adalah bentuk turunannya, yaitu arsitektur
\textit{plugin}. Eclipse, Visual Studio Code, Chrome, dan Firefox memakai bentuk
tersebut. Pada keempatnya, inti yang dipasang pengguna sebenarnya kecil, dan
hampir seluruh kemampuan yang dipakai sehari-hari datang dari bagian tambahan.

Pertanyaan yang dijawab bab ini karena itu bukan bagaimana menulis kernel,
melainkan bagaimana sebuah inti dapat menerima kemampuan baru tanpa mengenal
satu pun penyumbangnya. Pembahasan mengikuti kerangka yang disusun oleh Richards
dan Ford~\cite{richards2020fundamentals}, serta catatan Bolour tentang
arsitektur \textit{plugin} pada Eclipse~\cite{bolour2003eclipse}.

\section{Inti yang Kecil, \textit{Plugin} yang Banyak}

Susunannya terdiri atas dua bagian saja. \textbf{Inti} memuat kemampuan yang
berlaku bagi semua pemakai, misalnya menemukan bagian tambahan, menyimpan
daftarnya, dan menjalankan salah satunya atas permintaan. \textbf{\textit{Plugin}}
memuat kemampuan yang hanya dibutuhkan sebagian pemakai, misalnya dukungan satu
bahasa pemrograman tertentu.

Gambar~\ref{fig:inti-plugin} menunjukkan susunannya. Seluruh panah berangkat
dari \textit{plugin} menuju inti, dan tidak ada satu pun panah yang berangkat
dari inti menuju \textit{plugin}.

\begin{figure}[htbp]
  \centering
  \begin{tikzpicture}[
    kotak/.style={draw, rounded corners=2pt, minimum width=2.2cm,
      minimum height=0.7cm, align=center, font=\scriptsize\bfseries},
    panah/.style={-{Stealth[length=2mm]}, thick}
  ]
    \node[kotak, fill=praditagreen!25, minimum width=3.2cm,
      minimum height=1.1cm] (inti) {Inti\\
      {\tiny\mdseries menemukan, mendaftar, menjalankan}};

    \node[kotak, fill=blue!8, above left=1.1cm and 0.2cm of inti] (p1)
      {\textit{plugin} konversi};
    \node[kotak, fill=blue!8, above right=1.1cm and 0.2cm of inti] (p2)
      {\textit{plugin} kurs};
    \node[kotak, fill=blue!8, below left=1.1cm and 0.2cm of inti] (p3)
      {\textit{plugin} riwayat};
    \node[kotak, fill=green!10, below right=1.1cm and 0.2cm of inti] (p4)
      {\textit{plugin} baru};

    \draw[panah] (p1) -- (inti);
    \draw[panah] (p2) -- (inti);
    \draw[panah] (p3) -- (inti);
    \draw[panah] (p4) -- (inti);
  \end{tikzpicture}
  \caption{Inti beserta keempat \textit{plugin}. Seluruh panah menunjuk ke
    dalam, sehingga menambah kotak baru tidak menyentuh isi inti}
  \label{fig:inti-plugin}
\end{figure}

Analogi yang mendekati adalah stopkontak di dinding sebuah ruang kuliah. Dinding
menyediakan lubang dengan bentuk yang sudah disepakati, dan pemasang dinding
tidak perlu tahu perangkat apa yang kelak ditancapkan. Proyektor, kipas, dan
pengisi daya telepon sama-sama masuk selama bentuk stekernya sesuai. Dinding
juga tidak ikut rusak ketika satu perangkat terbakar, sebab pengaman memutus
arus ke lubang itu saja.

Pembagian tersebut menentukan letak setiap keputusan.
Tabel~\ref{tab:inti-plugin} merangkum tanggung jawab dan larangan kedua bagian.

\begin{table}[htbp]
  \centering
  \caption{Tanggung jawab dan larangan inti serta \textit{plugin}}
  \label{tab:inti-plugin}
  \begin{tabularx}{\textwidth}{|l|X|X|}
    \hline
    \textbf{Bagian} & \textbf{Tanggung jawab} & \textbf{Tidak boleh mengetahui} \\
    \hline
    Inti & Menemukan manifes, menyimpan registri, menjalankan perintah & Nama, jumlah, dan isi \textit{plugin} mana pun \\
    \hline
    \textit{Plugin} & Menyediakan satu kemampuan beserta manifesnya & Isi \textit{plugin} lain, dan cara inti menyimpan registrinya \\
    \hline
    Kontrak & Menyatakan bentuk yang wajib dipenuhi kedua belah pihak & Teknologi penyimpanan maupun bentuk antarmuka \\
    \hline
  \end{tabularx}
\end{table}

\section{Kontrak, Manifes, dan Registri}

Tiga hal dibutuhkan agar susunan tersebut bekerja. \textbf{Kontrak} menyatakan
bentuk yang wajib dipenuhi setiap \textit{plugin}. \textbf{Manifes} menyatakan
kemampuan apa yang disumbangkan, dan disimpan di luar kode. \textbf{Registri}
menyimpan manifes yang sudah terbaca, sehingga inti dapat mencari tanpa
memindai ulang.

Eclipse menyebut kontrak tersebut \textit{extension point}, yaitu slot yang
disediakan sebuah \textit{plugin} agar \textit{plugin} lain dapat
mengisinya~\cite{bolour2003eclipse}. Isi manifes dibaca lebih dahulu, lalu
tersedia lewat sebuah registri, dan cara itulah yang membuat inti dapat
menemukan penyumbang tanpa mengimpor kelasnya.

Listing~\ref{lst:kontrak} memperlihatkan kontrak pada contoh bab ini. Isinya
hanya satu bentuk pemanggilan, sebab kontrak yang gemuk menyulitkan penulis
\textit{plugin}.

\begin{lstlisting}[style=PythonStyle, caption={Kontrak yang wajib dipenuhi
  setiap \textit{plugin}, tersedia di \berkas{sources/chapter05/kontrak.py}},
  label=lst:kontrak]
# Kunci yang wajib ada pada setiap berkas manifes.
KUNCI_MANIFES = ("nama", "perintah", "keterangan", "modul", "kelas")


class Plugin(Protocol):
  """Satu kemampuan yang disumbangkan sebuah plugin kepada inti."""

  def __init__(self, inti: object) -> None:
    """Menerima inti sebagai satu-satunya ketergantungan sebuah plugin."""

  def jalankan(self, argumen: list[str]) -> str:
    """Menjalankan kemampuannya lalu mengembalikan teks siap tampil."""
\end{lstlisting}

Manifes sengaja disimpan sebagai berkas JavaScript Object Notation (JSON) yang
terpisah dari modulnya. Pemisahan itu bukan soal selera. Manifes yang terpisah
dapat dibaca tanpa menjalankan kode pluginnya, dan kemampuan itulah yang
dipakai pada seksi berikutnya. Visual Studio Code menyimpannya di
\texttt{package.json}, Chrome dan Firefox di
\texttt{manifest.json}~\cite{chrome2026content,mozilla2026manifest}, sedangkan
Eclipse dahulu memakai \texttt{plugin.xml}~\cite{bolour2003eclipse}.

Gambar~\ref{fig:kelas-microkernel} menunjukkan hubungan ketiganya sebagai
diagram kelas. Perhatikan bahwa panah dari inti hanya menuju kontrak dan
manifes, tidak pernah menuju modul \textit{plugin}.

\begin{figure}[htbp]
  \centering
  \includegraphics[width=0.78\textwidth, height=0.52\textheight,
    keepaspectratio]{kelas-microkernel.pdf}
  \caption{Inti mengenal kontrak dan manifes saja. Ketiga \textit{plugin}
    mengisi kontrak yang sama, dan manifeslah yang menunjuk nama modulnya}
  \label{fig:kelas-microkernel}
\end{figure}

\section{Kapan \textit{Plugin} Diaktifkan}

Menemukan \textit{plugin} dan mengaktifkannya adalah dua langkah yang berbeda,
dan jarak antara keduanya menentukan waktu siap sebuah aplikasi.
\textbf{Aktivasi awal} mengimpor seluruh modul saat aplikasi mulai.
\textbf{Aktivasi lambat} menunda impor sampai kemampuannya benar-benar dipakai.

Visual Studio Code memilih cara kedua. Setiap ekstensi menyatakan
\textit{activation event} pada manifesnya, misalnya \texttt{onLanguage} atau
\texttt{onCommand}, sehingga ekstensi baru dimuat ketika peristiwa itu terjadi.
Dokumentasinya bahkan menandai peristiwa yang menyala di segala keadaan sebagai
pilihan terakhir, sebab pilihan itu memperlambat awal
mula~\cite{vscode2026activation}. Eclipse menempuh jalan yang sama lewat
pemuatan yang ditunda, dan objek \textit{plugin} baru dibuat saat kemampuannya
dipakai~\cite{bolour2003eclipse}.

Mekanisme pemuatannya sudah lama dipakai di luar Python. Contoh Java pada
\berkas{sources/chapter05/java-example/} memuat setiap \textit{plugin} dari
berkas JAR tersendiri memakai \texttt{URLClassLoader}, membaca kunci
\texttt{mainclass} pada \texttt{plugin.properties}, lalu membentuk objeknya
lewat \texttt{Class.forName} beserta konstruktor hasil refleksi.
Tabel~\ref{tab:java-python} memetakan kedua sisinya.

\begin{table}[htbp]
  \centering
  \caption{Padanan mekanisme pemuatan pada contoh Java dan contoh Python}
  \label{tab:java-python}
  \begin{tabularx}{\textwidth}{|X|X|}
    \hline
    \textbf{Contoh Java} & \textbf{Contoh Python} \\
    \hline
    \texttt{plugin.properties} berisi kunci \texttt{mainclass} & \texttt{plugins/plugin\_*.json} berisi kunci \texttt{modul} dan \texttt{kelas} \\
    \hline
    \texttt{URLClassLoader} beserta \texttt{Class.forName} & \texttt{importlib.import\_module} beserta \texttt{getattr} \\
    \hline
    Konstruktor hasil refleksi menerima \texttt{Greeter} & Konstruktor kelas menerima objek \texttt{Inti} \\
    \hline
    Antarmuka \texttt{Plugin} & \textit{Protocol} \texttt{Plugin} pada \berkas{sources/chapter05/kontrak.py} \\
    \hline
  \end{tabularx}
\end{table}

Gambar~\ref{fig:sekuens-aktivasi} menunjukkan urutannya pada contoh bab ini.
Manifes dibaca saat mulai, sedangkan impor modulnya menunggu pemanggilan
pertama.

\begin{figure}[htbp]
  \centering
  \includegraphics[width=0.95\textwidth, height=0.55\textheight,
    keepaspectratio]{sekuens-aktivasi.pdf}
  \caption{Aktivasi lambat. Impor hanya terjadi pada pemanggilan pertama, dan
    \textit{plugin} yang tidak pernah dipanggil tidak pernah diimpor}
  \label{fig:sekuens-aktivasi}
\end{figure}

Listing~\ref{lst:aktifkan} memperlihatkan intinya. Hasil impor disimpan,
sehingga pemanggilan kedua tidak membayar biaya yang sama.

\begin{lstlisting}[style=PythonStyle, caption={Aktivasi satu \textit{plugin},
  tersedia di \berkas{sources/chapter05/inti.py}}, label=lst:aktifkan]
def aktifkan(self, perintah: str) -> kontrak.Plugin:
  """Mengimpor modul, mengambil kelasnya, lalu membuat objek pluginnya."""
  if perintah in self.plugin_aktif:
    return self.plugin_aktif[perintah]  # objek pemenuh kontrak.Plugin
  manifes = self.manifes[perintah]
  try:
    modul = importlib.import_module(manifes["modul"])
    kelas = getattr(modul, manifes["kelas"])
    plugin = kelas(self)
  except Exception as kesalahan:
    self.galat_muat[perintah] = str(kesalahan)
    raise PluginGagal(f"{manifes['nama']} gagal dimuat: {kesalahan}")
  self.plugin_aktif[perintah] = plugin
  return plugin
\end{lstlisting}

\section{Ketika Satu \textit{Plugin} Rusak}

Kemampuan menerima kode dari luar membawa persoalan yang tidak dimiliki
arsitektur pada bab sebelumnya. Kode itu ditulis pihak lain, kualitasnya tidak
dapat dipastikan, dan sebagiannya pasti akan gagal. Pertanyaannya, apakah
kegagalan satu \textit{plugin} ikut menjatuhkan intinya.

Jawabannya ditentukan oleh dua keputusan. Keputusan pertama adalah kapan modul
diimpor, sebab kegagalan saat impor jauh lebih berbahaya daripada kegagalan saat
pemanggilan. Keputusan kedua adalah apakah pemanggilannya dibungkus penangkap
galat. Chrome menempuh pemisahan yang lebih jauh lagi, yaitu menjalankan
\textit{content script} di dalam \textit{isolated world} yang terpisah dari
halaman maupun dari ekstensi lain~\cite{chrome2026content}.

Kasus berikut memperlihatkan akibatnya. Sebuah tim memasang lima \textit{plugin}
pada aplikasi internalnya, dan seluruhnya diaktifkan saat aplikasi mulai. Satu
\textit{plugin} membaca berkas pengaturan pada tingkat modul, di luar fungsi
mana pun. Pada suatu pagi berkas pengaturan itu terhapus, dan impornya gagal.
Aplikasi tidak menampilkan galat satu \textit{plugin}, melainkan gagal mulai
sama sekali. Empat \textit{plugin} lain yang sehat ikut tidak berjalan, dan
pengguna tidak punya cara untuk menonaktifkan yang rusak, sebab layarnya tidak
pernah muncul.

Tabel~\ref{tab:isolasi} merangkum kedua jalur tersebut beserta akibatnya.

\begin{table}[htbp]
  \centering
  \caption{Nasib aplikasi ketika satu \textit{plugin} gagal diimpor}
  \label{tab:isolasi}
  \begin{tabularx}{\textwidth}{|l|X|X|}
    \hline
    \textbf{Jalur} & \textbf{Yang terjadi} & \textbf{Akibat bagi pengguna} \\
    \hline
    Tanpa pembungkus & Galat naik sampai ke pemanggil paling luar & Aplikasi berhenti, \textit{plugin} sehat ikut mati \\
    \hline
    Dengan pembungkus & Galat dicatat, perintah lain tetap dilayani & Satu kemampuan hilang, sisanya tetap jalan \\
    \hline
    Proses terpisah & Galat tertahan di dalam prosesnya sendiri & Aplikasi utuh, tetapi biayanya lebih mahal \\
    \hline
  \end{tabularx}
\end{table}

\section{\textit{Plugin} pada Produk Nyata}

Keempat produk berikut memakai susunan yang sama dengan nama yang berbeda.
Tabel~\ref{tab:produk-plugin} merangkum istilah dan berkas manifes masing-masing.

\begin{table}[htbp]
  \centering
  \caption{Sebutan dan berkas manifes pada empat produk yang memakai
    \textit{plugin}}
  \label{tab:produk-plugin}
  \begin{tabularx}{\textwidth}{|l|l|X|}
    \hline
    \textbf{Produk} & \textbf{Manifes} & \textbf{Catatan} \\
    \hline
    Eclipse & \texttt{plugin.xml} & Kontraknya disebut \textit{extension point}, dan pengisinya disebut \textit{extension} \\
    \hline
    Visual Studio Code & \texttt{package.json} & Pemuatan ditunda sampai \textit{activation event} yang didaftarkan terjadi \\
    \hline
    Chrome & \texttt{manifest.json} & \textit{Content script} berjalan di \textit{isolated world} yang terpisah \\
    \hline
    Firefox & \texttt{manifest.json} & Memakai model WebExtensions yang sama, dengan sejumlah perbedaan yang disengaja \\
    \hline
  \end{tabularx}
\end{table}

Perbedaan namanya menutupi kesamaan susunannya. Ketiganya menyimpan manifes di
luar kode, menunda pemuatan, dan menolak menyebut nama penyumbang di dalam
intinya. Yang berbeda hanyalah seberapa jauh pemisahannya dibawa, dan Chrome
membawanya paling jauh dengan memisahkan lingkungan eksekusinya.

\section{Ringkasan}

Microkernel membagi sistem menjadi inti yang kecil beserta bagian tambahan yang
dapat dipasang dan dicabut. Bentuk yang paling sering ditemui bukan kernel
sistem operasi, melainkan arsitektur \textit{plugin} pada perkakas sehari-hari.
Eclipse, Visual Studio Code, Chrome, dan Firefox sama-sama memakainya.

Yang membuat susunan tersebut bekerja adalah arah pengenalan. Inti menulis
kontrak dan membaca manifes, sedangkan \textit{plugin} yang mengisi kontrak
tersebut. Karena inti tidak pernah menyebut nama satu penyumbang pun, menambah
kemampuan berarti menambah berkas, bukan menyunting inti. Latihan~1 membuktikan
klaim itu dengan menghitung penyebutannya, dan angkanya nol.

Manifes yang terpisah dari kode membuka pilihan kedua, yaitu kapan sebuah
\textit{plugin} diaktifkan. Aktivasi lambat membaca manifes saat mulai lalu
menunda impor sampai kemampuannya dipakai. Pada contoh bab ini, cara tersebut
membuat waktu siap 3,3 kali lebih cepat, dan selisihnya lebih besar daripada
\textit{range} pengukurannya sendiri. Kesimpulan itu berbeda dengan Bab~2, dan
perbedaannya justru memperlihatkan gunanya membandingkan selisih dengan
\textit{range}.

Kemampuan menerima kode dari luar menuntut harga tersendiri. Satu
\textit{plugin} yang gagal saat impor dapat menjatuhkan seluruh aplikasi bila
pemanggilannya tidak dibungkus. Latihan~2 memperlihatkan selisihnya, yaitu tiga
perintah tetap berjalan pada jalur berpelindung, sedangkan jalur tanpa
pelindung berhenti pada berkas yang rusak.

Pesan utama bab ini karena itu bukan bahwa inti harus kecil, melainkan bahwa
inti harus buta terhadap penyumbangnya. Kebutaan itu yang dibayar dengan
kontrak, manifes, dan registri, dan ketiganya baru terasa mahal ketika jumlah
\textit{plugin} masih sedikit.

\section{Latihan}

Ketiga latihan berikut memakai satu basis kode yang sama di
\berkas{sources/chapter05/}, tetapi menyasar tiga masalah yang berbeda.
Latihan~1 memeriksa penyebutan \textit{plugin} di dalam inti. Latihan~2
memeriksa akibat satu \textit{plugin} yang rusak. Latihan~3 mengukur selisih
waktu siap antara kedua cara aktivasi. Hasil setiap mahasiswa dapat berbeda,
dan yang dilaporkan adalah pola perbandingannya, bukan angka mutlaknya.

Venv yang sama dipakai kembali, dan pengaktifannya seperti pada
Listing~\ref{lst:venv-ch05}. Latihan bab ini hanya memakai pustaka standar
Python, sehingga tidak ada pemasangan tambahan.

\begin{lstlisting}[language=bash, caption={Mengaktifkan venv sebelum ketiga
  latihan, skrip tersedia di \berkas{sources/siapkan_venv.sh}},
  label=lst:venv-ch05]
bash sources/siapkan_venv.sh
source .venv/bin/activate
cd sources/chapter05
python inti.py daftar
# Keluarannya baris pertama: Plugin terpasang: 4
\end{lstlisting}

\subsection*{Latihan 1: Menghitung Penyebutan \textit{Plugin} di Dalam Inti}

Latihan ini menguji klaim utama arsitektur \textit{plugin}, yaitu menambah
kemampuan cukup dengan menambah berkas. Klaim tersebut menyebut angka, sehingga
dapat dibuktikan salah. Satu penyebutan nama modul \textit{plugin} di dalam inti
sudah cukup untuk menggugurkannya. Langkah kerjanya sebagai berikut.

\begin{enumerate}
  \item Jalankan \texttt{python periksa\_inti.py}, lalu catat baris
    \texttt{Plugin disebut inti} beserta ketiga angka di bawahnya.
  \item Buka \berkas{sources/chapter05/inti.py}, lalu cari nama modul
    \textit{plugin} di dalamnya. Catat hasil pencarian itu.
  \item Salin \berkas{sources/chapter05/plugins/plugin_riwayat.py} beserta
    manifesnya menjadi \texttt{plugin\_pajak.py}, lalu ubah nama, perintah, dan
    nama kelasnya.
  \item Jalankan \texttt{python inti.py daftar}, lalu catat jumlah
    \textit{plugin} yang terbaca.
  \item Jalankan ulang \texttt{python periksa\_inti.py}, lalu catat berapa baris
    inti yang berubah selama langkah ketiga dan keempat.
\end{enumerate}

Tabel~\ref{tab:lat1plugin-contoh} berisi hasil satu kali pemeriksaan nyata.

\begin{table}[htbp]
  \centering
  \caption{Contoh hasil Latihan 1 yang sudah terisi}
  \label{tab:lat1plugin-contoh}
  \begin{tabularx}{\textwidth}{|X|l|}
    \hline
    \textbf{Yang dilaporkan} & \textbf{Nilai} \\
    \hline
    Manifes ditemukan & 4 \\
    \hline
    \textit{Plugin} disebut inti & 0 \\
    \hline
    Baris inti & 128 \\
    \hline
    Baris seluruh \textit{plugin} & 104 \\
    \hline
    Bagian \textit{plugin} & 45 persen \\
    \hline
  \end{tabularx}
\end{table}

Isikan hasil pemeriksaan pada Tabel~\ref{tab:lat1plugin-hasil}.

\begin{table}[htbp]
  \centering
  \caption{Lembar isian Latihan 1}
  \label{tab:lat1plugin-hasil}
  \begin{tabularx}{\textwidth}{|X|l|l|}
    \hline
    \textbf{Yang dilaporkan} & \textbf{Sebelum} & \textbf{Sesudah} \\
    \hline
    Manifes ditemukan & \dots & \dots \\
    \hline
    \textit{Plugin} disebut inti & \dots & \dots \\
    \hline
    Baris inti & \dots & \dots \\
    \hline
    Baris seluruh \textit{plugin} & \dots & \dots \\
    \hline
  \end{tabularx}
\end{table}

Jawab tiga pertanyaan berikut dengan menunjuk angka pada tabel. Pertama, baris
inti tidak berubah sama sekali walaupun satu \textit{plugin} ditambahkan.
Sebutkan bagian inti yang membuat hal itu mungkin. Kedua, jelaskan apa yang akan
terjadi pada angka penyebutan bila inti mengimpor satu \textit{plugin} secara
langsung demi kemudahan. Ketiga, hitung bagian kode yang berupa \textit{plugin},
lalu bandingkan dengan perkiraan awal sebelum pengukuran dilakukan.

\subsection*{Latihan 2: Menjatuhkan Aplikasi dengan Satu \textit{Plugin}}

Latihan ini membuktikan bahwa kegagalan satu \textit{plugin} tidak harus
menjatuhkan intinya. Berkas \berkas{sources/chapter05/plugin_rusak.py} sengaja
gagal pada tingkat modul, sehingga kegagalannya muncul saat impor. Langkah
kerjanya sebagai berikut.

\begin{enumerate}
  \item Baca \berkas{sources/chapter05/plugins/plugin_rusak.py}, lalu
    tunjukkan baris yang menyebabkan kegagalannya.
  \item Jalankan \texttt{python uji\_isolasi.py}, lalu catat kedua baris
    hasilnya.
  \item Jalankan \texttt{python inti.py rusak} dan \texttt{python inti.py kurs}
    berurutan, lalu catat apakah perintah kedua tetap dilayani.
  \item Buka \berkas{sources/chapter05/inti.py}, lalu temukan bagian yang
    membungkus galat impor.
  \item Hapus sementara pembungkus tersebut, jalankan ulang
    \texttt{python inti.py kurs}, lalu catat akibatnya. Kembalikan kodenya
    sesudah dicatat.
\end{enumerate}

Tabel~\ref{tab:lat2plugin-contoh} berisi hasil satu kali percobaan nyata.

\begin{table}[htbp]
  \centering
  \caption{Contoh hasil Latihan 2 yang sudah terisi}
  \label{tab:lat2plugin-contoh}
  \begin{tabularx}{\textwidth}{|l|l|X|}
    \hline
    \textbf{Jalur} & \textbf{Perintah jalan} & \textbf{Catatan} \\
    \hline
    Dengan pelindung & 3 dari 4 & Satu \textit{plugin} dilewati, sisanya dilayani \\
    \hline
    Tanpa pelindung & 3 modul & Berhenti pada \texttt{plugin\_rusak} \\
    \hline
  \end{tabularx}
\end{table}

Isikan hasil percobaan pada Tabel~\ref{tab:lat2plugin-hasil}.

\begin{table}[htbp]
  \centering
  \caption{Lembar isian Latihan 2}
  \label{tab:lat2plugin-hasil}
  \begin{tabularx}{\textwidth}{|l|l|X|}
    \hline
    \textbf{Jalur} & \textbf{Perintah jalan} & \textbf{Catatan} \\
    \hline
    Dengan pelindung & \dots & \dots \\
    \hline
    Tanpa pelindung & \dots & \dots \\
    \hline
    Tanpa pembungkus galat & \dots & \dots \\
    \hline
  \end{tabularx}
\end{table}

Jawab tiga pertanyaan berikut. Pertama, jelaskan mengapa kegagalan saat impor
lebih berbahaya daripada kegagalan saat pemanggilan. Kedua, sebutkan apa yang
hilang bagi pengguna ketika satu \textit{plugin} dilewati, lalu bandingkan
dengan yang hilang ketika aplikasi tidak mau mulai. Ketiga, jelaskan mengapa
pemisahan proses seperti pada Chrome lebih kuat daripada pembungkus galat,
beserta harga yang dibayarnya.

\subsection*{Latihan 3: Mengukur Biaya Aktivasi Awal}

Latihan ini mengukur selisih waktu siap antara aktivasi lambat dan aktivasi
awal. Pengukurannya diulang seratus putaran, sebab satu putaran tidak cukup
untuk menyimpulkan apa pun. Cache impor Python dibersihkan pada setiap putaran,
agar kedua jalur menghadapi keadaan yang sama. Langkah kerjanya sebagai berikut.

\begin{enumerate}
  \item Jalankan \texttt{python ukur\_aktivasi.py}. Skrip tersebut mencetak
    baris kemajuan setiap sepuluh putaran, sehingga jalannya terlihat.
  \item Catat kedelapan angka pada laporan akhirnya.
  \item Buka \berkas{sources/chapter05/ukur_aktivasi.py}, lalu tunjukkan baris
    yang membersihkan cache impor.
  \item Hapus sementara pembersihan tersebut, jalankan ulang, lalu catat
    perubahan angkanya. Kembalikan kodenya sesudah dicatat.
\end{enumerate}

Tabel~\ref{tab:lat3plugin-contoh} berisi hasil satu kali pengukuran nyata.

\begin{table}[htbp]
  \centering
  \caption{Contoh hasil Latihan 3 yang sudah terisi, seratus putaran}
  \label{tab:lat3plugin-contoh}
  \begin{tabularx}{\textwidth}{|X|l|}
    \hline
    \textbf{Yang dilaporkan} & \textbf{Nilai} \\
    \hline
    Median aktivasi lambat & 0,4006 ms \\
    \hline
    Median aktivasi awal & 1,3406 ms \\
    \hline
    Selisih median antar jalur & 0,9399 ms \\
    \hline
    Aktivasi awal lebih lambat & 3,3 kali \\
    \hline
    \textit{Range} di dalam jalur lambat & 0,2896 ms \\
    \hline
    \textit{Range} dibagi selisih & 0,3 kali \\
    \hline
    Putaran lambat lebih cepat & 100 dari 100 \\
    \hline
  \end{tabularx}
\end{table}

Baris keenam itulah yang membedakan bab ini dari Bab~2. Di sana \textit{range}
di dalam satu jalur puluhan kali lebih besar daripada selisih yang dicari,
sehingga selisihnya belum dapat disebut ada. Di sini \textit{range}-nya justru
sepertiga selisihnya, dan pencacahan per putaran menegaskannya dengan angka
100 dari 100. Selisih yang seperti inilah yang layak dilaporkan.

Isikan hasil pengukuran pada Tabel~\ref{tab:lat3plugin-hasil}.

\begin{table}[htbp]
  \centering
  \caption{Lembar isian Latihan 3}
  \label{tab:lat3plugin-hasil}
  \begin{tabularx}{\textwidth}{|X|l|l|}
    \hline
    \textbf{Yang dilaporkan} & \textbf{Dengan pembersihan} & \textbf{Tanpa pembersihan} \\
    \hline
    Median aktivasi lambat & \dots & \dots \\
    \hline
    Median aktivasi awal & \dots & \dots \\
    \hline
    Selisih median antar jalur & \dots & \dots \\
    \hline
    \textit{Range} dibagi selisih & \dots & \dots \\
    \hline
    Putaran lambat lebih cepat & \dots & \dots \\
    \hline
  \end{tabularx}
\end{table}

Jawab tiga pertanyaan berikut dengan menunjuk angka pada tabel. Pertama,
jelaskan mengapa selisih pada bab ini layak disebut ada, sedangkan selisih pada
Bab~2 tidak. Kedua, jelaskan apa yang terjadi pada angka ketika pembersihan
cache dihapus, lalu sebutkan pengukuran mana yang menjadi tidak jujur karenanya.
Ketiga, dengan menggabungkan hasil ketiga latihan, tentukan kapan aktivasi awal
tetap layak dipilih walaupun lebih lambat, lalu sebutkan angka yang menjadi
dasar keputusannya.
